% Source: https://artofproblemsolving.com/community/c3h1848454_some_problems
% Stable identity: aops-topic-1848454-q05
Repeatedly take the integer part and invert the remainder, rationalizing each time. Write each complete quotient as $x_k=\dfrac{\sqrt{43}+m}{d}$; its integer part $a_k$ is found from $\sqrt{43}\approx6.557$.
\[
\begin{aligned}
x_{0}&=\sqrt{43}, &\qquad a_{0}&=6\\[6pt]
x_{1}&=\dfrac{1}{\sqrt{43}-6}=\dfrac{\sqrt{43}+6}{7}, &\qquad a_{1}&=1\\[6pt]
x_{2}&=\dfrac{7}{\sqrt{43}-1}=\dfrac{\sqrt{43}+1}{6}, &\qquad a_{2}&=1\\[6pt]
x_{3}&=\dfrac{6}{\sqrt{43}-5}=\dfrac{\sqrt{43}+5}{3}, &\qquad a_{3}&=3\\[6pt]
x_{4}&=\dfrac{3}{\sqrt{43}-4}=\dfrac{\sqrt{43}+4}{9}, &\qquad a_{4}&=1\\[6pt]
x_{5}&=\dfrac{9}{\sqrt{43}-5}=\dfrac{\sqrt{43}+5}{2}, &\qquad a_{5}&=5\\[6pt]
x_{6}&=\dfrac{2}{\sqrt{43}-5}=\dfrac{\sqrt{43}+5}{9}, &\qquad a_{6}&=1\\[6pt]
x_{7}&=\dfrac{9}{\sqrt{43}-4}=\dfrac{\sqrt{43}+4}{3}, &\qquad a_{7}&=3\\[6pt]
x_{8}&=\dfrac{3}{\sqrt{43}-5}=\dfrac{\sqrt{43}+5}{6}, &\qquad a_{8}&=1\\[6pt]
x_{9}&=\dfrac{6}{\sqrt{43}-1}=\dfrac{\sqrt{43}+1}{7}, &\qquad a_{9}&=1\\[6pt]
x_{10}&=\dfrac{7}{\sqrt{43}-6}=\sqrt{43}+6, &\qquad a_{10}&=12.
\end{aligned}
\]
For example, in row $1$: $\dfrac{1}{\sqrt{43}-6}=\dfrac{\sqrt{43}+6}{43-36}=\dfrac{\sqrt{43}+6}{7}\approx1.79$, and the remainder $x_1-1=\dfrac{\sqrt{43}-1}{7}$ is inverted in row $2$. In row $10$ the remainder is $x_{10}-12=\sqrt{43}-6$, the same as in row $0$, so from here on the quotients repeat. Hence
\[
\sqrt{43}=\boxed{[6;\ \overline{1,1,3,1,5,1,3,1,1,12}]},
\]
with period $10$. (Check: the convergent $[6;1,1,3,1,5,1,3,1,1]=\dfrac{3482}{531}$ satisfies $3482^2-43\cdot531^2=1$.)
